\scanpage{21}
\diagram{\begin{tikzcd}[ampersand replacement=\&,row sep=large,column sep=huge]
 H^*(X,\mathbf C) \arrow[r,"\sim"] \& H^*(X,\Omega_X^\bullet) \\
 \underline{\underline H}^*(Y,i^!(\mathbf C)) \arrow[u,"\unsure{?}"'] \arrow[r,"\sim","{\text{comm (\unsure{trivial})}}"'] \& \underline{\underline H}^*(Y,i^!(\Omega_X^\bullet)) \arrow[u] \\
 H^*_{\deleted{\unsure{Y}}}(\unsure{Y},\mathbf C) \arrow[u,"P_{\mathrm{top}}"] \arrow[r,"\sim","{\frac{1}{2i\pi}\mathrm{can}}"'] \& H^*(Y,\Omega_Y^\bullet) \arrow[u,"P_{\mathrm{an}}"]
\end{tikzcd}}

qui provient de la comm. de
\diagram{\begin{tikzcd}[ampersand replacement=\&,row sep=large,column sep=huge]
 i^!(\mathbf C) \arrow[r,"\mathrm{can}","\sim"'] \& i^!(\Omega_X^\bullet) \\
 \mathbf C_Y \arrow[u,"P_{\mathrm{top}}"] \arrow[r,dash,"\sim","{\left(\frac{1}{2i\pi}\right)^p\mathrm{can}}"'] \& \Omega_Y^\bullet \arrow[u]
\end{tikzcd}}
lui $\unsure{=}$ contenu dans la formule fondamentale

\uline{N.B.} On voit que l'\uline{image directe analytique} ne \uline{\unsure{transforme pas}} \unsure{classes} entières en \uline{\unsure{classes}} entières : il faut un facteur $(2i\pi)^p$ \unsure{en|en $+$} pour qu'il en soit ainsi.

\diagram{\fbox{$\begin{tikzcd}[ampersand replacement=\&,row sep=large,column sep=large]
 \text{\unsure{cycles analytiques de codim }i} \arrow[r,dash] \& A^i(X) \arrow[r,"P_{\mathrm{top}}^{\mathbf Z}"] \arrow[d,"{\left(\frac{1}{2i\pi}\right)^pP_{\mathrm{an}}}"] \& H^{2i}(X,\mathbf Z) \arrow[d,"\mathrm{can}"] \& \text{est comm.} \\
 \& H^{\unsure{2i}}(X,\Omega_X^\bullet) \arrow[r,"\sim"] \& H^{2i}(X,\mathbf C)
\end{tikzcd}$}}

\mspage{20}
